% !PRINTED 338
\weberpage{338}
%<<CONT
in dem $f(x)$ und $f''(x)$ \gesperrt{dasselbe Vorzeichen} haben. Die
beiden Fig. 23 und 24 veranschaulichen das Verh\"altniss bei positivem
 $f'(x)$ und positivem und negativem $f''(x)$.

Es ist also im ersten Falle $\beta'$, im zweiten $\alpha'$ ein besserer
Ann\"aherungswerth, als $\beta$ und $\alpha$.

Will man gleichzeitig die andere Grenze verschieben, um
ein neues engeres Intervall zu bekommen, so kann man nach
Fourier von dem anderen Endpunkte, also im Falle der Fig. 23
in $a$ die zu $b\beta'$ parallele gerade Linie ziehen, und erh\"alt den
Punkt $\alpha'$ als unteren Grenzpunkt des neuen Intervalles. Es ist
dann im ersten Falle
\[
\alpha' = \alpha - \frac{f(\alpha)}{f'(\beta)}, \quad \beta' = \beta - \frac{f(\beta)}{f'(\beta)},
\]
im zweiten Falle
\[
\alpha' = \alpha - \frac{f(\alpha)}{f'(\alpha)}, \quad \beta' = \beta - \frac{f(\beta)}{f'(\alpha)},
\]
und $(\alpha', \beta')$ ist ein engeres Intervall, in dem die gesuchte Wurzel liegt.

% SOURCE FIGURE CROP: Bd. I registered scan 364 (printed p. 338), 600 dpi,
% crop x=250,y=2250,w=1200,h=1150; source-faithful raster, Fig. 25.
\weberfigureimage{weber-b1-fig25.png}{0.42}

Es ist kaum n\"othig, die beiden
anderen F\"alle, in denen $f'(x)$ im Intervall
 negativ ist, noch besonders zu
betrachten.

Man kann aber auch, um zwei
neue Grenzen zu erhalten, mit noch
besserem Erfolge die Newton'sche
Methode mit der Interpolationsmethode
verbinden, wie die Fig. 25 zeigt.

Man erh\"alt dann als die beiden
neuen Grenzen:
\[
\begin{gathered}
\alpha' = \alpha - \frac{f(\alpha)(\beta - \alpha)}{f(\beta) - f(\alpha)} = \frac{\alpha f(\beta) - \beta f(\alpha)}{f(\beta) - f(\alpha)}, \\
\beta' = \beta - \frac{f(\beta)}{f'(\beta)}.
\end{gathered}
\]

Wir fassen die Resultate in folgendem Satze zusammen:

\gesperrt{Ist ein Intervall $(\alpha, \beta)$ abgegrenzt, in dem
$f(x)$ einmal, $f'(x)$ und $f''(x)$ gar nicht ihr Zeichen
wechseln, und ist $\beta$ der Endpunkt des Intervalles
, in dem $f(\beta)$ und $f''(\beta)$ dasselbe Vorzeichen
haben, gleichviel, ob $\beta$ kleiner oder gr\"osser als}
%CONT>>
